\chapter{Experimental Methodology}
\label{sec:methodology}

Our simulator, application workloads, simulation results, and results processing scripts are publicly available~\cite{HIDDEN_URL}.

%are publicly available~\footnote{URL omitted for double-blind review.}.



\section{Application Workloads}

Some of our results are illustrated for the fusion
science application in Fig.~\ref{fig:use_case}. Tasks
$T_1$, $T_2$, and $T_3$ have four, four, and three execution modes, respectively,
each described by affine $t$, $d$, and $e$ functions.  These particular
execution modes and affine functions are based on results in~\cite{9912722}
and on information provided to us by the authors.  
%Although we have
%used some best guesses for some coefficients in these functions, we
%believe our specification to be representative of a practical use case.

We also evaluate heuristics over a large set of synthetic application
configurations.  A synthetic configuration is generated as follows. We
randomly sample the number of tasks $M{\in}\{2,3,4,5\}$. For each task
$T_i$ we randomly sample its number of modes $K_i{\in}\{2, 3, 4\}$. We then
generate a first ``safe" ($s$) execution mode (fast, high error)
and a last ``risky" ($r$) execution mode (slow, low error) as
follows, omitting arguments on which a function does not depend:

\begin{minipage}{0.39\columnwidth}
\begin{itemize}[leftmargin=0pt, labelindent=0pt, labelsep=0.3em]
\item $t^i_s{=}u_{50, 600}$
\item $d^i_s{=}\lceil u_{10^{10}, 10^{12}} \rceil$
\item $e^i_s(y){=}u_{3, 10} \cdot y$
\end{itemize}
\end{minipage}
\begin{minipage}{0.55\columnwidth}
\begin{itemize}[leftmargin=0pt, labelindent=0pt, labelsep=0.3em]
	\item $t^i_r{=}u_{2 \cdot t^i_s , 6 \cdot t^i_s}$
	\item $d^i_r{=}\lceil u_{0.01 \cdot d^i_s, 0.3 \cdot d^i_s} \rceil$
	\item $e^i_r(y){=}u_{1, \min(0.6 \cdot e^i_s(1), e^i_s(1) - 0.5)} \cdot y$
\end{itemize}
\end{minipage}~\\

\noindent
where $u_{a,b}$ denotes a value sampled once from the uniform distribution 
on $[a,b] \subset \mathbb{R}$.  Additional modes are generated between the safe and risky
modes using linear equidistant interpolation. In practice, the $t$ and $d$
functions are likely non-constant (as in our fusion science use
case). However, using constant functions, especially in the input size,
makes it straightforward to generate a broad range of configurations
in which no execution mode is trivially dominated by another.

The intervals for the uniform distribution are selected empirically
to ensure that the application configurations are diverse enough that
different scheduling algorithms can make a wide range of reasonable
decisions.  For each generated application configuration we
compute $p_s$, resp. $p_r$, the probability that the execution on a single
PE succeeds by the deadline when using only safe, resp. risky, execution
modes.  $p_s$ and $p_r$ are computed using dynamic programming (see
Chapter~\ref{sec:onehost}) using the same $D$, $\lambda$, $R$, $B_r$, and
$B_w$ values as that used in simulation.  To eliminate trivially easy or
hard cases, we filter out configurations with $p_s{<}0.15$, $p_s{>}0.85$,
or $p_r{<}0.01$.  The $D$ and $\lambda$ values merely establish a common
time scale; failure behavior is driven by the overall application
execution time relative to the deadline, $D$, and to the mean time between failures,
$1/\lambda$.  Our synthetic configurations do not aim to model any
specific workload. Instead, they define a set of scenarios with different
relative execution-time, data-size, and error magnitudes,  over which
different scheduling principles can be compared systematically.

%\TODO{Perhaps (likely cannot due to space) add data to justify that the configurations are really diverse....}

\section{Simulation Methodology}

Our simulator builds on a simulation framework designed to model real-world
software/hardware platforms, and in particular network and I/O operations,
accurately.  Scheduling works in the literature, like this one, typically
use an abstract model of I/O operations, where transfer time is simulated as
the data volume divided by some nominal bottleneck bandwidth, to render
analytical developments and algorithm design
tractable~\cite{10.1145/2501654.2501664, 7965082}. Our simulator supports
two I/O simulation modes.  Its ``model-consistent" mode enforces the above
abstract model, and is used to validate our implementation. Its
``realistic" mode uses realistic simulation models that account for
protocol and contention effects.  Unless specified otherwise, simulation
results are obtained with the realistic mode.  Each invocation of the
simulator takes in a full specification of the scheduling problem, as well
as a seed value for the random number generator used to inject failures
into the execution.  For a given instance of the scheduling problem,
simulation results are obtained for different seeds, i.e., different
failure traces.

Unless specified otherwise, analytical/simulation results are 
for $D{=}10^4$, $\lambda{=}10^{-3}$, $R{=}0$, $B_r{=}$2 GB/s, and $B_w{=}$1 GB/s.

